%==============================================================================
\section{Introduction}\label{sec:intro}
%==============================================================================

\subsection{The problem}\label{ssec:intro-problem}

A workshop has a batch of parts to make. Each part needs its own handful of cutting
tools: a drill of one diameter, a particular milling cutter, a tap for a thread. The
machine can use any tool the workshop owns, but not all of them at once. Only the tools
loaded into its magazine --- a rack or carousel with a fixed number of slots --- are
within reach. Reaching any other one means stopping the machine and swapping a loaded
tool out for it.

That swap is where the problem comes from. The machine produces nothing while it
happens, and a batch of parts with varied requirements forces a great many of them. The
standard model of the situation therefore counts swaps and nothing else, treating
cutting times and part handling as fixed because they do not depend on the order the
parts are made in. The seven modelling assumptions behind that choice are due to
\citet{Crama1994} and are set out in Section~\ref{ssec:ssp}.

Two decisions control the count. The first is the order the parts are made in. The
second is which tools to keep in the magazine and which to displace at each step. The
Job Sequencing and Tool Switching Problem, written SSP throughout this report, makes
both decisions together and minimises the total number of tool insertions.
Figure~\ref{fig:setup} shows it at its smallest.

\begin{figure}[ht]
\centering
\begin{tikzpicture}[
   font=\small, >=Stealth,
   job/.style={draw=cBlue!70, fill=white, rounded corners, semithick,
               minimum width=2.7cm, minimum height=1.0cm, align=center},
   chip/.style={draw=black!25, rounded corners=1pt, minimum size=0.72cm,
                inner sep=0pt, font=\bfseries, text=white},
   swaparr/.style={->, line width=1.1pt, cRed}]
  \node[job] (j1) at (0,2)   {\textbf{job 1} needs\\ {\color{cRed}$a$}, {\color{cBlue}$b$}};
  \node[job] (j2) at (5,2)   {\textbf{job 2} needs\\ {\color{cBlue}$b$}, {\color{cGreen}$c$}};
  \node[job] (j3) at (10,2)  {\textbf{job 3} needs\\ {\color{cGreen}$c$}, {\color{cOrange}$d$}};
  \node[chip,fill=cRed]    (m1a) at (-0.42,0) {$a$};
  \node[chip,fill=cBlue]   (m1b) at ( 0.42,0) {$b$};
  \node[chip,fill=cBlue]   (m2a) at ( 4.58,0) {$b$};
  \node[chip,fill=cGreen]  (m2b) at ( 5.42,0) {$c$};
  \node[chip,fill=cGreen]  (m3a) at ( 9.58,0) {$c$};
  \node[chip,fill=cOrange] (m3b) at (10.42,0) {$d$};
  \node[draw=black!45, rounded corners, fit=(m1a)(m1b), inner sep=2.5pt] (m1) {};
  \node[draw=black!45, rounded corners, fit=(m2a)(m2b), inner sep=2.5pt] (m2) {};
  \node[draw=black!45, rounded corners, fit=(m3a)(m3b), inner sep=2.5pt] (m3) {};
  \draw[->,black!45] (j1) -- (m1);  \draw[->,black!45] (j2) -- (m2);  \draw[->,black!45] (j3) -- (m3);
  \draw[swaparr] (m1) -- (m2)
     node[midway,above,font=\scriptsize,align=center]{swap\\ $-a$, $+c$};
  \draw[swaparr] (m2) -- (m3)
     node[midway,above,font=\scriptsize,align=center]{swap\\ $-b$, $+d$};
  \node[font=\scriptsize,text=black!55] at (5,-1.05) {the magazine (grey box) holds only $b=2$ tools at a time};
\end{tikzpicture}
\caption{The problem at its smallest. Three jobs run from left to right on a
machine whose magazine holds two tools. A job can run only when both of its tools
are loaded. Filling the magazine for job~1 costs two insertions. Each later job
happens to share one tool with the job before it, so one tool is swapped in at each
step: two more insertions, four in total. Changing the order of the jobs, or making
the magazine larger, changes that count. Finding the order that minimises the total,
over all $n!$ orders, is the SSP.}
\label{fig:setup}
\end{figure}

Figure~\ref{fig:twoorders} takes that same instance and changes only the order, which
changes the count.

\begin{figure}[ht]
\centering
\begin{tikzpicture}[
   font=\small, >=Stealth,
   slot/.style={draw=black!45, rounded corners=1pt, minimum size=0.62cm,
                inner sep=0pt, font=\small},
   fresh/.style={slot, fill=black!72, text=white, draw=black!72},
   held/.style={slot, fill=black!8},
   jl/.style={font=\scriptsize\bfseries, text=black!75}]

 % ---------------- row 1 : order 1,2,3 -> cost 4 ----------------
 \begin{scope}
   \node[font=\small\bfseries, anchor=west] at (-2.35,1.02) {order $1,2,3$};
   \node[anchor=west, font=\small] at (2.55,1.02) {four tool insertions \textbf{(optimal)}};
   \foreach \x/\lab in {0/{job 1},2.3/{job 2},4.6/{job 3}}
     \node[jl] at (\x+0.31,0.62) {\lab};
   % job1 {a,b}: both new
   \node[fresh] (r1a) at (0,0)    {$a$};  \node[fresh] (r1b) at (0.62,0) {$b$};
   % job2 {b,c}: b held, c new
   \node[held]  (r2b) at (2.3,0)  {$b$};  \node[fresh] (r2c) at (2.92,0) {$c$};
   % job3 {c,d}: c held, d new
   \node[held]  (r3c) at (4.6,0)  {$c$};  \node[fresh] (r3d) at (5.22,0) {$d$};
   \foreach \a/\b in {r1b/r2b, r2c/r3c}
     \draw[->, black!45, shorten >=2pt, shorten <=2pt] (\a) to[bend left=38] (\b);
   \node[font=\scriptsize, text=black!60, anchor=west] at (6.15,0) {$2+1+1=4$};
 \end{scope}

 % ---------------- row 2 : order 1,3,2 -> cost 5 ----------------
 \begin{scope}[yshift=-2.15cm]
   \node[font=\small\bfseries, anchor=west] at (-2.35,1.02) {order $1,3,2$};
   \node[anchor=west, font=\small] at (2.55,1.02) {five --- one insertion wasted};
   \foreach \x/\lab in {0/{job 1},2.3/{job 3},4.6/{job 2}}
     \node[jl] at (\x+0.31,0.62) {\lab};
   \node[fresh] (s1a) at (0,0)    {$a$};  \node[fresh] (s1b) at (0.62,0) {$b$};
   \node[fresh] (s2c) at (2.3,0)  {$c$};  \node[fresh] (s2d) at (2.92,0) {$d$};
   \node[held]  (s3c) at (4.6,0)  {$c$};  \node[fresh] (s3b) at (5.22,0) {$b$};
   \draw[->, black!45, shorten >=2pt, shorten <=2pt] (s2c) to[bend left=38] (s3c);
   \node[font=\scriptsize, text=black!60, anchor=west] at (6.15,0) {$2+2+1=5$};
 \end{scope}

 % ---------------- legend ----------------
 \node[fresh] (lg1) at (-1.6,-3.35) {};
 \node[anchor=west, font=\scriptsize, text=black!65] at (-1.25,-3.35) {inserted here (counted)};
 \node[held]  (lg2) at (2.75,-3.35) {};
 \node[anchor=west, font=\scriptsize, text=black!65] at (3.1,-3.35) {already loaded (free)};
\end{tikzpicture}
\caption{The same three jobs, two orders. Running job~2 between jobs~1 and~3 lets the
magazine hand tool~$b$ forward and then tool~$c$ forward, so only one tool is fetched at
each step. Running job~3 second breaks both chains: the magazine must be emptied and
refilled, and tool~$b$ --- discarded after job~1 --- has to be fetched a second time.
That repeat is the only difference between the two rows, and it is what an exact method
has to reason about. There are $n!$ orders and no way to see which chains survive
without searching.}
\label{fig:twoorders}
\end{figure}

Figure~\ref{fig:landscape} does the same on a real benchmark instance of eight jobs,
small enough that all $40{,}320$ orders can be enumerated. It shows how much lies
between a good order and a bad one, and where the bound that every exact method starts
from sits relative to both.

\begin{figure}[ht]\centering
\begin{tikzpicture}
\begin{axis}[reportaxis, height=5.9cm,
  ybar, bar width=8.5pt,
  xlabel={tool switches required by an ordering},
  ylabel={number of orderings},
  xmin=8.6, xmax=25, ymin=0, ymax=10600,
  xtick={10,13,14,16,18,20,22,24},
  ytick={0,2000,4000,6000,8000},
  scaled y ticks=false,
  yticklabel style={/pgf/number format/fixed,
                    /pgf/number format/1000 sep={,}},
  clip=false]
  \addplot[draw=black!65, fill=black!22] table[x=cost,y=count] {figdata/landscape.dat};
  % the coverage bound: below everything achievable
  \draw[black, line width=1pt, dashed] (axis cs:10,0) -- (axis cs:10,9700);
  \node[anchor=west, font=\scriptsize, align=left] at (axis cs:10.25,9100)
    {\textbf{coverage bound $q=10$}\\[-1pt] what every compact\\[-1pt] relaxation returns};
  % the optimum
  \draw[black, line width=1pt] (axis cs:13,0) -- (axis cs:13,6100);
  \node[anchor=north east, font=\scriptsize, align=right] at (axis cs:12.7,6100)
    {\textbf{optimum $\Zst=13$}\\[-1pt] reached by $180$\\[-1pt] of the $40{,}320$};
  % the median ordering
  \draw[black!55, line width=0.8pt, densely dotted] (axis cs:18,0) -- (axis cs:18,9700);
  \node[anchor=west, font=\scriptsize, text=black!60, align=left] at (axis cs:18.3,9250)
    {median ordering: $18$};
\end{axis}
\end{tikzpicture}
\caption{Every one of the $40{,}320$ orderings of instance \texttt{L1-1}
($n=8$ jobs, $\lvert\U\rvert=15$ tools, magazine $b=5$), enumerated, and the tool
switches each requires. Three things are visible at once. A randomly chosen ordering
costs $18$, and the best costs $13$, so the choice of order is worth about a third of
the total. Only $180$ orderings out of forty thousand are optimal, which is why the
outer problem is hard. And the counting bound of $q=10$, which is what the SSPMF
compact baseline returns, lies \emph{below the entire
distribution}: no ordering whatsoever attains it. A solver given that bound cannot
prove any ordering optimal until it has closed a gap of three by branching alone.
Section~\ref{sec:comp} shows what that costs in practice and Section~\ref{sec:disc} why
the bound behaves this way.}
\label{fig:landscape}
\end{figure}

The same structure appears well outside machining, and everything in this report
carries over. In printed-circuit-board assembly the tools are component feeders and
the jobs are the boards. In colour printing they are ink cartridges on a press, and
the problem occurs there in exactly the form studied here
\citep{Burger2015}. In warehouse picking they are the items a picker can carry at
once. What matters in each case is that a limited, reusable resource serves groups
of demands, and that the order of the demands can be chosen.

\subsection{Why the problem is hard, and why it is still open}\label{ssec:intro-why}

The two decisions behave very differently.

Loading is easy. Once the order of the jobs is fixed, the cheapest loading policy is
known in closed form. It is the Keep Tool Needed Soonest rule of \citet{Tang1988}:
when a tool must be inserted into a full magazine, remove the tool whose next use
lies furthest in the future. The rule is optimal, and it runs in
$O(n\lvert T\rvert)$ time.

Ordering is hard. \citet{Crama1994} proved the SSP NP-hard for every fixed magazine
capacity $b\ge2$. The difficulty is that a good order places jobs with similar tool
requirements next to one another, but ``similar'' here means nothing more than
overlapping sets, with no natural distance to measure it, and the magazine limit
couples jobs that are far apart in the order.

That combination --- an easy inner problem inside a hard outer one --- is exactly
what decomposition methods are built for, and it is the reason the SSP has attracted
exact methods for four decades. Two gaps in that literature motivate this work.

The first is on the practical side. Exact methods for the integrated SSP exist and
are used as baselines in Section~\ref{sec:comp}, but practice does not generally rely on
them. The standard
approach groups the jobs into as few magazine-fitting batches as possible, then
orders the batches. It is fast, and where it has been measured against exact optima in
industrial printing it does well \citep{Burger2015}. The leading hybrid genetic search
instead starts from random job orders improved by local search \citep{Mecler2021}; it
does not use grouping solutions as seeds. \citet{CramaKlundert1999} analyse the
group-then-sequence construction and prove
that it is a $b$-approximation. What has not been done is to sharpen that ratio, to say
on which instances the method is exact, or to distinguish the two different values the
method can be read as returning --- and the third of these turns out to change the first
two. How far
from optimal it can be, and on which instances it is exact, are open questions.

The second is on the exact side. The three compact integer formulations reimplemented
as campaign baselines have weak linear relaxations, and the strongest of that baseline
set, SSPMF \citep{daSilva2024}, relaxes exactly to the number of tools minus the
magazine capacity, under its assumption that every tool is used. That is an
elementary counting bound: every tool that does not fit in the initial magazine must be
inserted at some point. Thus the best relaxation among the three reimplemented
baselines is this counting argument, and a solver relying on it cannot prove optimality
where the bound is loose. This is deliberately not a literature-wide ranking:
the symmetry-strengthened JGSMF model of \citet{Akhundov2024}, for example, was not
implemented in this internship.

\subsection{What was done during this internship}\label{ssec:intro-done}

The internship took place at the Laboratoire d'Informatique, de Mod\'elisation et
d'Optimisation des Syst\`emes (LIMOS), under the advisorship of Prof.~Rafael
Colares and Prof.~Annegret Wagler, with Dr.~Renaud Chicoisne as academic
supervisor. The research directions were set by the advisors, and this report is
organised around them. They are recorded here explicitly, so that the origin of
each line of work is clear.

\begin{itemize}[topsep=4pt,itemsep=3pt]
  \item \textbf{The configuration-based view of the problem} was proposed by
  Prof.~Colares and Prof.~Wagler, and frames the whole project. A
  \emph{configuration} is a full magazine, and a schedule is a walk through
  configurations that serves every job. Section~\ref{ssec:gtsp} develops it.

  \item \textbf{The reduction of the SSP to a generalised travelling-salesman
  problem} was set by Prof.~Colares. It is the formal content of that view, and it
  is what makes both exact methods below natural.

  \item \textbf{The analysis of the gap between the two-phase heuristic and the
  optimum} was set by Prof.~Wagler, together with the request to enumerate the
  integral solutions of the grouping problem with PORTA. Her specific questions
  were to construct worst-case examples with the number of groups fixed, to study
  instances where the gap is attained, and to prove an upper bound tighter than the
  trivial $b(\Kst-1)$. Section~\ref{sec:gap} answers the first two and gives
  several answers to the third.

  \item \textbf{The study of the grouping problem through clutters and blocking
  duality} was also set by Prof.~Wagler. Section~\ref{ssec:clutter} reports it.

  \item \textbf{The branch-and-Benders-cut algorithm} was proposed by
  Prof.~Colares, following the internship report of \citet{Kashou2025}, which
  applies the same architecture to a production-line balancing problem. The
  suggestion was to try lazy cuts first and cuts at fractional solutions second.
  Section~\ref{ssec:bbc} develops it and Section~\ref{ssec:results-bbc} measures it.

  \item \textbf{The position-indexed formulations} were proposed by
  Prof.~Colares, including the idea that makes them work: fixing a row set that
  is polynomial in the instance size, so that column generation only ever adds
  columns to rows that already exist. The purpose was to make branch-and-price over
  configurations possible. Section~\ref{ssec:pos} develops them.
\end{itemize}

\noindent
Work on these directions fell into three parts, and each has a section of its own
below.

The first was theoretical: the structural analysis of Section~\ref{sec:gap}. It
answers the questions Prof.~Wagler set, and raises several that emerged from them.

The second was implementation. The branch-and-Benders-cut solver, the two
branch-and-price prototypes, three formulations from the literature reimplemented as
baselines, a hybrid genetic heuristic, the campaign harness and the cluster pipeline
were built and validated against brute-force optima before any of them was used to
produce a result. Section~\ref{sec:exact} describes the methods and
Appendix~\ref{app:repro} the code.

The third was experimental: the benchmark campaign of Section~\ref{sec:comp}, run on
the LIMOS cluster over the standard instance families under a uniform protocol,
together with the small-scale enumeration used throughout Section~\ref{sec:gap} to
test each statement as it was made. One correction made during the campaign is worth
recording, because it changed the results substantially: the Benders master was found
to be missing the tool coverage bound, and the campaign was re-run once it was added.

\subsection{Conventions}\label{ssec:intro-read}

Section~\ref{sec:prelim} defines the problem and everything used to attack it,
including a self-contained account of the integer-programming background in
Section~\ref{ssec:toolkit}, which a reader who already knows integer programming may
skip. Section~\ref{sec:sota} reviews the literature. Longer proofs are in
Appendix~\ref{app:proofs}, the verification program in Appendix~\ref{app:verify}, and
Appendix~\ref{app:repro} says how to reproduce every experiment.

Every formal result carries a label --- Theorem, Proposition, Lemma, Corollary,
Observation, or Conjecture --- and the labels are used strictly. A Proposition has a
proof. An Observation was established by running a program. A Conjecture has neither,
only evidence. Where an argument is given but has not been independently checked it is
headed \emph{Supporting argument} rather than \emph{Proof}. The distinction is kept
even where it makes the text less confident.
